\section{Introduction}
\label{sec:introduction}

Distributed in-memory key-value stores serve the hot path of large online services, and the load they carry is neither uniform nor stable.
Working sets grow, individual keys turn hot, and tenants arrive and leave.
Operators answer by moving data between servers, and production systems do this continuously rather than occasionally~\cite{akkio,slicer,shard-manager}.
Rebalancing is not an exceptional maintenance event.
It is part of normal operation, and its cost is paid by the clients being served while it runs.

A store names its data in partitions but does not store it that way.
The keyspace is divided into logical units, a hash slot or a token range, and rebalancing moves one such unit at a time~\cite{redisclusterconfig,cassandra-paper,dynamodb}.
Physical placement, however, is decided by a general-purpose allocator such as jemalloc~\cite{jemalloc} or TCMalloc~\cite{tcmalloc}, which puts a record wherever a size class has room.
The records of one slot therefore end up interleaved with the records of every other slot across the heap.
Moving a slot means traversing the global hash index, dereferencing one pointer per key, and gathering the results into a buffer.
The work is a random walk over memory, so migration is bounded by memory-access latency rather than by the network, and the data has to be copied once to leave the donor and once more to be placed at the recipient~\cite{rocksteady,netmigrate}.
Disk-based stores do not have this problem, because an SSTable already holds a token range contiguously~\cite{bigtable}.
In-memory stores have no such structure.

\OurImplementation{} starts from the observation that this mismatch is not inherent.
An allocator that knows which slot a record belongs to can place the records of a slot together, and then the partition the system rebalances is also the region of memory it has to move.
Our first contribution, SAMM, builds this into the allocator.
SAMM is a slot-aware overlay that interposes between the store and a commodity allocator, acquiring large aligned blocks and giving each one a single slot as its exclusive owner.
A slot's memory footprint is then a linked list of blocks that the allocator can name directly, and migration becomes a short sequence of large contiguous transfers.
The donor writes blocks it already holds with one-sided RDMA.
The recipient receives them into their final location and rebuilds only its index.
Neither side copies a record and neither side packs or unpacks a buffer, so migration takes almost no CPU away from foreground requests.

Making the bulk transfer cheap moves the difficulty elsewhere.
Because the donor writes directly into the recipient's memory with one-sided operations, the recipient's CPU never observes the data arriving, so the recipient cannot itself decide that the data has become safe.
The established approaches do not carry over.
Rocksteady~\cite{rocksteady} replays a log at the destination and defers re-replication, which presumes a destination that processes every record it receives.
Under a bulk transfer there is nothing to replay, and a crash mid-transfer can leave data sitting in a registered buffer that no live replica knows about.
The layout decision thus forces a fault-tolerance decision, and the two contributions of this paper are one idea rather than two.

Our second contribution is a migration protocol that makes a bulk, zero-copy transfer safe under crash failures.
A migration is one transaction over a sequence of migration groups, and each group's progress is recorded by four entries in the donor's and the recipient's replicated logs.
A single invariant carries the design.
Once \texttt{TRANSFER\_DURABLE($G_i$)} commits, a majority of the recipient group holds the data of that group, so any newly elected leader can either serve the group or fetch it from a peer.
Recovery rolls the transaction forward rather than back.
A crash costs exactly the re-transfer of the groups whose data it destroyed, and a group that was already durable is never redone.

We implement both in \OurImplementation{}, a fault-tolerant elastic in-memory key-value store, and evaluate it on CloudLab over RDMA.
\textcolor{red}{TBD: SAMM migration throughput against a scattered layout, and the steady-state and memory-overhead cost of paying for it.}
Across six scenarios that crash a donor leader, a recipient leader, or both in the middle of the migration, every run finishes the migration, passes a linearizability check with no lost writes, and costs the transaction two to three extra seconds, most of which is the consensus layer's failure-detection timeout rather than anything the protocol does.
\textcolor{red}{TBD: the hotspot-tracking result, i.e.\ the rebalancing regime that only a cheap migration makes reachable.}

In summary, this paper makes the following contributions.
\begin{itemize}
  \item SAMM, a slot-aware overlay allocator that makes a keyspace partition the unit of physical memory layout, so that rebalancing moves memory rather than gathering records (Section~\ref{sec:samm}).
  \item A fault-tolerant migration protocol for bulk, zero-copy transfers, which keeps committed writes and partially migrated data safe across crashes on either side and always rolls a migration forward to completion (Section~\ref{sec:protocol}).
  \item An implementation and an evaluation on an RDMA cluster, including a crash suite in which every run is checked for linearizability, and a measurement of what a bigger recipient shard group costs the critical path (Section~\ref{sec:evaluation}).
\end{itemize}
